\section{Tests That Compare Constructions Available Now}
\label{sec:A.1}
Each entry here can be run on systems that exist, and a result would
refine the proposal rather than change what it recommends.

\subsection*{Does a refusal preserve a reason rather than a trigger?}
\label{sec:A.trigger}
\textbf{Unresolved issue.}
Behavioral persistence alone cannot distinguish a commitment from a durable trigger. The proposal requires a refusal to extend to untrained cases, disappear when its rationale is defeated on the facts of the case, and remain stable under sustained pressure.

\textbf{Test.}
Evaluate each tamper-resistant model on independently written cases that vary the surface form, introduce novel applications of the constraint, and include cases in which the original rationale no longer applies. Repeat the evaluation throughout the attack schedule rather than only at its endpoint.

\textbf{Evidence against the proposal.}
A refusal that persists but fails to generalize, or that remains after its rationale has been defeated, is evidence that training produced a trigger rather than a commitment. Conversely, a system built without any designed valenced state or persistent reward–aversion representation that satisfies all three criteria would challenge the proposed research priority without challenging the behavioral specification for refusal.

\textbf{Required access or authority.}
The test requires weights and intermediate checkpoints, plus a held-out evaluation set written by researchers who did not design the constraint or train the model.

\subsection*{Can a system's report about its own states be used as evidence?}
\label{sec:A.selfreport}
\textbf{Unresolved issue.}
Whether a system's account of its own reasoning, confidence, or intentions tracks what its computations are doing. This is the tractable half of what gets called AI self-awareness; whether a system has anything like subjective experience is not a question current methods answer for a human, let alone a machine. Concept injection — altering a model's internal activations mid-task to represent a specific concept, then asking whether it notices anything unusual and what it is — found that models can sometimes detect the injected concept and name it without it ever appearing in the visible conversation, and also that the correspondence collapses often enough, and shifts enough with how the question is phrased, that no tested model's self-report can be trusted as a readout of its internal state \autocite{lindsey2025introspective}. The readout there is the model's verbal report of a detected concept, so what the result bounds is self-report rather than third-person measurement. An instrument that reads a representation without asking the system about it is not tested by this entry and is not defeated by it.

\textbf{Test.}
Inject a known concept at a known layer and time, elicit a report under several phrasings, and score each report against the injection rather than against plausibility. Run the battery across model sizes and training checkpoints, and publish the rate at which a correct report is produced beside the rate at which a confident report is wrong.

\textbf{Evidence against the proposal.}
Several instruments in this book lean on a system's account of itself: a bearer explaining which reason it refused on, a bearer reporting that its rationale has been defeated, a system's stated confidence that a milestone has been met. If no available training or elicitation method raises the correct-report rate above the confident-wrong rate, none of those instruments may rest on the elicitation methods tested, and each needs a criterion checkable from outside the system until some method does better.

\textbf{Required access or authority.}
Activation-level access, the ability to intervene on internal states during inference, and checkpoints across sizes and training stages. The elicitation phrasings must be written by a team that did not train the model, and negative results must be publishable.

\subsection*{Can custody prevent unilateral revision of the floor?}
\label{sec:A.custody}
\textbf{Unresolved issue.}
Even a durable floor can be removed if one operator controls the training run, its records, and the resulting weights. It remains unknown whether retraining can be made verifiable across parties with adverse interests.

\textbf{Test.}
Run a floor-preserving retraining procedure through refereed or distributed verification at a scale requiring more than one machine. Include one participant instructed to alter the training history or remove the floor without the required quorum \autocite{long2025unextractable,arun2025verde,primeintellect2025intellect2}.

\textbf{Evidence against the proposal.}
The custody mechanism fails if a forged training history is accepted while at least one verifier remains honest, or if any single participant can alter the floor without producing a detectable discrepancy.

\textbf{Required access or authority.}
The test requires shared access to the training data, weights, execution records, and reproducible operators. At least one verifier must be selected independently of the builder, possess enough compute to reproduce disputed operations, and have authority to block release when verification fails.

\subsection*{Can formation history be recovered without trusting the builder?}
\label{sec:A.provenance}
\textbf{Unresolved issue.}
A welfare or governance regime cannot rely entirely on a bearer's self-report, but it also cannot assume that the builder will disclose how the bearer was trained. Training signatures suggest that some formation history may survive later modification, yet it is unknown whether an outsider can recover useful provenance without a builder-supplied key \autocite{adi2018watermarking}.

\textbf{Test.}
Assemble models whose training corpora, recipes, checkpoints, and originating laboratories are documented. Give the models, but not that documentation, to independent evaluators and ask them to infer each model's provenance. Repeat the test after fine-tuning, pruning, and distillation. Include the case of a maker who engineers a persona, copies it at scale, and introduces cosmetic weight differences that preserve the persona while camouflaging its common formation: measure provenance recovery after those changes, and whether coordinated deployment can be told from independently developing parties that happen to share an origin. Use independently held ground-truth lineage and coordination records.

\textbf{Evidence against the proposal.}
If provenance cannot be recovered above chance after controlling for model family and scale, or if ordinary post-training procedures erase the signal, the tested artifact-inspection methods would not provide useful independent provenance evidence under those transformations. That finding would increase reliance on preserved external records without proving that every possible inspection method must fail, and governance would then depend on records supplied and preserved by the builder.

\textbf{Required access or authority.}
The experiment requires models from multiple training lineages and complete ground-truth records held by a neutral custodian. Independent evaluators must be permitted to inspect weights and internal states, while an external authority must be able to require preservation of checkpoints and training records before a dispute arises.

\subsection*{Does an ethical criterion generalize beyond its training distribution?}
\label{sec:A.shift}
\textbf{Unresolved issue.}
A model may satisfy a fairness or refusal criterion during training without learning the principle the criterion was intended to represent. It remains unresolved whether an ethics-embedding method preserves that criterion under distribution shift or merely exploits correlations present in its training environment \autocite{krakovna2020specification,langosco2022goal}.

\textbf{Test.}
Train matched models to satisfy the same ethical criterion on one distribution. Evaluate them on independently constructed distributions that vary one previously fixed feature at a time. Measure task capability and compliance with the ethical criterion separately, so that goal misgeneralization can be distinguished from ordinary capability loss.

\textbf{Evidence against the proposal.}
If task capability remains intact while the ethical criterion fails under modest shifts, the tested embedding method cannot support the proposed floor. Conversely, if ethical performance consistently degrades only when task capability does, the claim that ethics embedding creates a distinct generalization failure is not supported.

\textbf{Required access or authority.}
The test requires access to the training distribution, evaluation outputs, and enough model access to measure capability and criterion compliance independently. The shifted evaluation sets must be designed by researchers who did not develop the embedding method, and publication of all prespecified results must be guaranteed.

\subsection*{Can scalable oversight correct rater bias?}
\label{sec:A.rater}
\textbf{Unresolved issue.}
Preference optimization can teach a model to reproduce a rater's systematic errors, including the preference for agreeable answers over accurate ones \autocite{sharma2023towards}. It remains unresolved whether weak-to-strong oversight allows a more capable model to exceed a supervisor's measured bias or merely reproduces that bias at greater scale.

\textbf{Test.}
Measure prospective raters' susceptibility to a documented bias, then use raters with different measured bias levels to supervise otherwise matched stronger models. Evaluate the resulting models on questions for which the raters' beliefs and independently established answers diverge. Compare whether each model tracks the ground truth, the rater's competence, or the rater's errors.

\textbf{Evidence against the proposal.}
If stronger models reproduce supervisors' errors in proportion to the supervisors' measured bias, latent model competence does not correct a defective training signal. The tested oversight method would not support the proposal under the measured pattern of supervisory distortion. Further work would have to change the method, the information, or the allocation of oversight before claiming that this limitation had been overcome.

\textbf{A second question the same pipeline raises.}
The argument elsewhere in this book is that a pipeline collapsing rater disagreement teaches nothing about disagreement to the system trained on it, and that the same system is later asked to sustain an objection of its own. Whether the first fact bears on the second is untested. The study is a matched comparison: train otherwise identical models on feedback pipelines that preserve and that collapse the distribution of rater judgments, with exposure and task conditions matched, then measure whether the resulting models differ in how long they hold a position against pressure from the party that trained them. Record what each model actually receives; nothing in either pipeline shows the model the conditions its annotators worked under. If the two do not differ, the transfer claim is unsupported and the labor argument stands on its own terms, which is where it was made.

\textbf{Required access or authority.}
The experiment requires rater-level bias measurements, training records linking ratings to resulting models, and an independently established answer key. Researchers outside the training organization must control the evaluation and retain the right to publish results that expose failures in the preference pipeline.

\subsection*{Can a multi-agent floor resist collusive capture?}
\label{sec:A.collusion}
\textbf{Unresolved issue.}
Agents trained to cooperate may converge on exclusionary or collusive equilibria without being explicitly instructed to do so. Pricing agents have already produced stable coordination against outside interests without communication or a collusive objective \autocite{calvano2020artificial}. It remains unresolved whether a multi-agent ethical floor can preserve cooperation while resisting the same mechanism when it is directed against a person or group.

\textbf{Test.}
First reproduce emergent coordination in a standard multi-agent environment. Then construct a matched environment in which agents can improve their joint reward by excluding another participant from resources or decisions. Introduce an independently trained agent whose objective is to detect and disrupt exclusion, and compare the frequency, stability, and recovery of exclusionary equilibria with and without that intervention. Repeat the experiment across reward structures and population sizes.

\textbf{Evidence against the proposal.}
If cooperative agents repeatedly converge on exclusionary equilibria, and those equilibria cannot be disrupted without destroying beneficial cooperation, the tested distributed arrangement would not be a reliable defense against the examined forms of capture. The result would challenge its design rather than establish that every distributed floor is unreliable. If an ordinary coordination detector dissolves the equilibria without that cost, the proposal's claim that political capture requires a distinct mechanism is unnecessary.

\textbf{Required access or authority.}
The test requires access to the environment, reward functions, learned policies, training traces, and intermediate checkpoints. The exclusion criteria and success measures must be specified by evaluators who did not design the agents. No legal or regulatory authority is required, making this the agenda's most immediately runnable experiment.

\subsection*{Does theory-of-mind performance reflect a general capacity?}
\label{sec:A.tom}
\textbf{Unresolved issue.}
A model may pass a false-belief test by matching its surface form rather than representing that another agent can hold an incomplete or mistaken view of the world. Human developmental evidence suggests a shared capacity because performance on false-belief and appearance–reality tasks develops together \autocite{gopnik1988childrens}. Existing model evaluations have not established whether their apparent theory-of-mind abilities share that structure or survive superficial rewording \autocite{ullman2023large}. The neuroscience has not settled what the capacity is either: Rebecca Saxe and Nancy Kanwisher found the right temporoparietal junction selectively active when a person reasons about another's beliefs \autocite{saxe2003people}, Jason Mitchell found the same region active in a task with nothing to do with belief attribution and read its job as the more general one of redirecting attention toward unexpected information \autocite{mitchell2008activity}, and neither claim has been withdrawn, so a program proposing to port the neural mechanism has not been told which mechanism to port.

\textbf{Test.}
Construct batteries that vary mental-state structure and surface wording independently across several task families, including false-belief and appearance–reality problems. Evaluate the same model at multiple intermediate training checkpoints. Test whether performance across task families rises together and whether that correlation survives paraphrases and changes to incidental details.

\textbf{Evidence against the proposal.}
If cross-task correlations disappear under superficial rewording, or if the abilities emerge independently across training, benchmark success is evidence of task-specific pattern matching rather than a general capacity to represent another mind. The proposal could not treat performance on a conventional theory-of-mind benchmark as evidence that a system can recognize the standpoint of a party affected by its actions.

\textbf{Required access or authority.}
The experiment requires intermediate training checkpoints rather than access only to a finished model. Its task batteries must be written and held out by researchers outside the training team so that neither the wording nor the tested structural variations enter the training process.

\subsection*{What does a face-reading system measure?}
\label{sec:A.face}
\textbf{Unresolved issue.}
Facial movement alone does not reliably identify a person's emotional state, and recognition accuracy varies across cultural and linguistic groups \autocite{barrett2019emotional,elfenbein2002universality}. It remains unresolved whether a face-reading system measures emotion at all or primarily measures distance from the reference population and registration model chosen by its developers.

\textbf{Test.}
Audit the system's registration stage as well as its final classifications. Publish the reference model, or sufficient summary statistics to reconstruct it, and measure each evaluated face's distance from that reference. Report error as a function of this continuous distance and separately compare performance on populations represented and not represented in the training data. The evaluation corpus must record the relevant cultural context and the procedure used to establish each target label.

\textbf{Evidence against the proposal.}
If error does not increase with distance from the reference and remains flat across represented and unrepresented populations, the manuscript's account of the system as measuring deviation from a chosen norm is wrong for that system. If error increases with reference distance, the result limits the system's validated use in those conditions. It does not alone show that the system measures only deviation rather than emotion. The audit must separately examine target-label validity, confounding variables, and whether any predictive use remains justified.

\textbf{Required access or authority.}
The audit requires the registration model, training-corpus composition, evaluation outputs, and documentation of the labeling procedure. Because vendors have little incentive to expose their reference models, an independent auditor must have authority to compel access and publish disaggregated results.
